% Physical pages 3 and 4 of the concise study: the movement of the formulary,
% stated as an argument and not as a summary of what follows. The developed
% comparison of the three readings begins on the fifth page.

\section{The Propers: Themes and Movement}
\zlabel{triptych:concise:themes:start}

This Mass confesses at its entrance a will that no one resists, and then shows
what lives inside that will: an adversary who must ask before he strikes, a
debtor forgiven beyond measure who must forgive his brother, and a people
brought out of Egypt that asks when its persecutors will be judged. It begins
with a creed spoken under sentence of death and ends with the food of
immortality, and between the two it teaches the Church who its enemy is, what it
owes, and where it is going.

The Introit is cut from the prayer that Mardochai prayed when the king's decree,
obtained by Aman, had condemned the Jews of the Persian empire to destruction
(Esth 13:9--11). In the Book of Esther he confesses that all things are in God's
power and that none can resist his will, ``if thou determine to save Israel''.
The Missal keeps the confession and drops the condition: \latin{In voluntate
tua, Domine, universa sunt posita, et non est qui possit resistere voluntati
tuae}. What Mardochai said of God's power to save one threatened people the
antiphon says of his will without condition, and it says it in the mouth of a
man whose enemy seems to have prevailed. In lines the Mass does not sing he
explains why he would not bow to Aman: ``I feared lest I should transfer the
honour of my God to a man'' (13:12--14). Its verse, the first of Ps~118, names
the road on which that confession is lived: ``Blessed are the undefiled in the
way, who walk in the law of the Lord.'' The Collect then calls the Church God's
family, \latin{familiam tuam}, and asks two things of his \latin{continua
pietas}: that under his protection, \latin{te protegente}, the family be free
\latin{a cunctis adversitatibus}, and that it be devoted to his name in good
works. A people that has confessed God's sovereignty asks at once to be guarded
within it, and the Collect's \latin{te protegente} and its adversities look
ahead to the Gradual's refuge and the Epistle's adversaries.

The Epistle names the adversity. St Paul, writing as a prisoner, ends his
letter to the Ephesians with a call to arms: ``Be strengthened in the Lord'';
put on ``the armour of God'' against the devil's deceits, \latin{insidias}, ``for
our wrestling is not against flesh and blood; but against principalities and
powers, against the rulers of the world of this darkness''. The Christian is to
resist ``in the evil day'' and, having done all, to stand, girt with truth, with
the breastplate of justice, the shoes of the gospel of peace, the shield of
faith, the helmet of salvation and the sword of the Spirit, ``which is the word
of God''. The strength is \latin{in Domino} and the armour is God's: both are
lent to the soldier. The Gradual answers with a fortress. ``Lord, thou hast
been our refuge from generation to generation'', it sings, and its second verse
sets God before the mountains, \latin{a saeculo et usque in saeculum tu es,
Deus}, in the present tense. The psalm goes on, in verses the Gradual leaves
aside, to the shortness of human life; what is sung is only the refuge and the
eternity. The Alleluia then remembers the first deliverance from an oppressor:
\latin{In exitu Israel de Aegypto}, the house of Jacob brought out from a
barbarous people.

At the centre stands the parable of the unforgiving servant (Mt 18:23--35). The
Mass begins it without Peter's question how often he must forgive, so that it
stands alone as a word of Jesus to his disciples. A king \latin{voluit rationem
ponere}, willed to settle accounts. A servant who owes ten thousand talents, a
sum no servant could repay, is ordered to be sold with his wife and children; he
falls down and asks for time, ``Have patience with me, and I will pay thee
all'', and is given remission of the whole. Going out, he meets a fellow servant
who owes him a hundred pence, takes him by the throat, hears the very words he
had spoken, and casts him into prison. The king calls him back, ``Thou wicked
servant, I forgave thee all the debt, because thou besoughtest me'', and hands
him to the torturers \latin{quoadusque redderet universum debitum}, until he
paid all the debt. The sentence then turns on every hearer: ``So also shall my
heavenly Father do to you, if you forgive not every one his brother from your
hearts.'' The whole debt that was forgiven is exacted again, and the parable
leaves its hearers to ask who revoked the pardon.

The Offertory returns to the adversary, telling the first chapters of Job in
words of its own. A man of the land of Hus, \latin{simplex et rectus, ac timens
Deum}; \latin{quem Satan petiit, ut tentaret: et data est ei potestas a Domino}
over his possessions and over his flesh; he lost his substance and his sons and
was struck with a grievous ulcer. The antiphon gives Satan a verb the book does
not use, \latin{petiit}, he asked, and it stops before Job's answer, ``the Lord
gave, and the Lord hath taken away \dots\ blessed be the name of the Lord''
(Job 1:21). Set
beside the Introit, it ties the hardest knot in the Mass: a will that none
resists, and an adversary who asks for a just man's ruin and is given power by
the Lord. The Secret offers the sacrifices by which God willed to be appeased,
\latin{quibus et te placari voluisti}, and asks him to restore salvation
\latin{potenti pietate}. The \latin{pietas} of the Collect returns, and it is
now called powerful.

The Communion returns to Ps~118, the psalm whose first verse the Introit sang,
and takes three half-verses from its strophe \latin{Caph}, whose speaker is worn out with waiting. It leaves out his fainting, his failing
eyes and the question ``How many are the days of thy servant?'', and sings a
hope, a question and a cry: the soul has hoped in God's salvation and in his
word; \latin{quando facies de persequentibus me iudicium?}; the wicked have
persecuted me, help me, O Lord my God. Sung after a Gospel that has commanded
forgiveness from the heart, the plea for judgment on persecutors would seem to
undo it. The Postcommunion closes on what has been received,
\latin{immortalitatis alimonia}, the food of immortality, and asks that what the
mouth has taken the pure mind may follow, \latin{pura mente sectemur}. Its pairing of mouth and mind is that of a
prayer the priest says after his Communion in the Order of Mass, \latin{Quod ore
sumpsimus} (\work{Ordo Missae}, no.~1130).

Three questions rise from this sequence, and each gives one reading of the
whole formulary, developed in the full study with four senses of its own.

The first, the war God permits and bounds, asks who the Christian's enemy is and how
the adversary's work can stand within a will that none resists. It is heard from
the Introit and the Offertory with the Epistle. St Gregory the Great answers
from Job that God looses the enemy and bridles him, permitting for his servant's
proving what the devil asks for his ruin; St John Chrysostom, from the Epistle,
that the spirits who work in persecutors are the real enemy, so that the wronged
man is to turn his anger on the devil and not on his brother.

The second, the
king's reckoning and the forgiven debtor, asks what the king's mercy requires of
the forgiven and whether a pardon once given can return. It is heard from the
Gospel. Chrysostom answers that the servant revoked the gift himself by his
cruelty; St Augustine, that God calls back what he remitted from the one who
will not forgive, and that forgiveness from the heart leaves room for correction
given in love.

The third, the Church's exodus toward the judgment, asks when the
judgment the Communion asks for will come, and where that leaves the people
that sings this Mass. It is heard from the Communion with the Alleluia and the
Postcommunion. Augustine hears in the plea the martyrs' cry under the altar,
answered with long-suffering until their brethren are complete; Chrysostom reads
the Epistle's girded loins as the girding of the Passover eaten on the night of
the Exodus, and the Christian life as a march toward departure.

The three hold much in common. All read the formulary under one sovereign God
who is refuge: the Introit's will that none resists and the Gradual's God before
the mountains stand behind each of them. All hear the Gospel's mercy and the
Communion's plea together, and none lets the plea undo the command to forgive.
They differ in the element each hears first and in the question each answers:
the first a question about God's government, the second a question about the
forgiven man's heart, the third a question about time and its end. The first and
second meet at a sentence of Chrysostom's on the enemy, the second and third at a
sermon of Augustine's on anger and the martyrs not yet avenged. Their answers
differ too, and
at one point the difference lies between the witnesses themselves: who the
Communion's persecutors are, and when the judgment it asks for will come. The
commentary that follows sets the readings' answers beside one another where the
texts put the same question to all three.
\zlabel{triptych:concise:themes:end}
